\documentclass[11pt,a4paper]{article}
\usepackage[utf8]{inputenc}
\usepackage[T1]{fontenc}
\usepackage[spanish]{babel}
\usepackage{lmodern}
\usepackage{geometry}
\usepackage{amsmath,amssymb,booktabs,array,longtable,xcolor,hyperref,enumitem,fancyhdr}
\geometry{margin=2.3cm}
\definecolor{azul}{HTML}{476885}
\hypersetup{colorlinks=true,linkcolor=azul,urlcolor=azul}
\pagestyle{fancy}\fancyhf{}\lhead{Teoría de Probabilidades}\rhead{Simulación en Excel}\cfoot{\thepage}
\newcommand{\excel}[1]{\par\smallskip\noindent\fcolorbox{azul}{azul!5}{\parbox{0.94\linewidth}{\ttfamily\small\detokenize{#1}}}\smallskip}
\title{\textcolor{azul}{Guía de aprendizaje}\\Simulación de una fila de servicio en Excel}
\author{Curso de Teoría de Probabilidades}
\date{}
\begin{document}\maketitle
\noindent\textbf{Actividad:} Simulación de 20 clientes. \quad \textbf{Modelo:} un servidor, una fila FIFO.\\
\textbf{Herramienta:} Microsoft Excel en español. \quad \textbf{Unidad:} minutos.
\section{Objetivo}
Construir una simulación de eventos discretos de una fila de servicio. Los tiempos \emph{entre llegadas} siguen una distribución exponencial y los tiempos de servicio se generan mediante una distribución normal. Se calculan las horas de llegada, las esperas y el tiempo total de cada cliente en el sistema.
\section{Supuestos y variables}
\begin{center}\begin{tabular}{ll}\toprule Parámetro & Valor inicial\\\midrule Número de clientes & 20\\ Tiempo medio entre llegadas & 3 minutos\\ Media del tiempo de servicio & 2,5 minutos\\ Desviación estándar del servicio & 0,5 minutos\\ Número de servidores & 1\\ Disciplina de atención & FIFO\\\bottomrule\end{tabular}\end{center}
Sea $X_i\sim\operatorname{Exp}(\lambda=1/3)$ el tiempo entre las llegadas de los clientes $i-1$ e $i$. El tiempo de llegada acumulado es $A_i=\sum_{j=1}^{i}X_j$. El tiempo de servicio se genera desde $S_i^*\sim N(2.5,0.5^2)$ y se usa $S_i=\max(0,S_i^*)$ para evitar duraciones negativas. Esta última transformación no es una normal estrictamente truncada: asigna a cero la probabilidad de valores negativos, que aquí es muy pequeña.

El inicio de servicio $B_i$, la espera $W_i$, la salida $D_i$ y el tiempo total $T_i$ se calculan como
\[
B_1=A_1,\qquad B_i=\max(A_i,D_{i-1})\quad(i\ge2),
\]
\[
W_i=B_i-A_i,\qquad D_i=B_i+S_i,\qquad T_i=D_i-A_i=W_i+S_i.
\]
\section{Estructura de la hoja de Excel}
En la fila 1 escriba estos encabezados:
\begin{center}\small\begin{tabular}{cl}\toprule Columna & Contenido\\\midrule A & Cliente\\ B & Tiempo entre llegadas\\ C & Tiempo de llegada acumulado\\ D & Tiempo de servicio\\ E & Inicio de atención\\ F & Tiempo de espera\\ G & Fin de atención\\ H & Tiempo total en el sistema\\\bottomrule\end{tabular}\end{center}
La fila 2 corresponde al cliente 1 y la fila 21 al cliente 20. En \texttt{A2} escriba 1 y en \texttt{A3} escriba \texttt{=A2+1}; copie hasta \texttt{A21}.
\section{Procedimiento paso a paso}
\subsection*{Paso 1. Tiempo entre llegadas: columna B}
En \texttt{B2} escriba:
\excel{=-3*LN(1-ALEATORIO())}
Copie hasta \texttt{B21}. La transformación inversa $-3\ln(1-U)$, con $U\sim\operatorname{Uniforme}(0,1)$, genera tiempos exponenciales con media 3 minutos.
\subsection*{Paso 2. Llegada acumulada: columna C}
En \texttt{C2} escriba:
\excel{=B2}
En \texttt{C3} escriba:
\excel{=C2+B3}
Copie \texttt{C3} hasta \texttt{C21}. La llegada es un instante medido desde el inicio de la observación, no el intervalo entre llegadas.
\subsection*{Paso 3. Tiempo de servicio: columna D}
En \texttt{D2} escriba:
\excel{=MAX(0;INV.NORM(ALEATORIO();2,5;0,5))}
Copie hasta \texttt{D21}. La función \texttt{INV.NORM} convierte un uniforme en una variable normal con media 2,5 y desviación 0,5. El \texttt{MAX} evita valores negativos. Dependiendo de la versión de Excel, el nombre de la función puede variar.
\subsection*{Paso 4. Inicio de atención: columna E}
En \texttt{E2} escriba:
\excel{=C2}
En \texttt{E3} escriba:
\excel{=MAX(C3;G2)}
Copie hasta \texttt{E21}. El servidor atiende al cliente cuando este llega o cuando termina el servicio anterior, lo que ocurra después.
\subsection*{Paso 5. Tiempo de espera: columna F}
En \texttt{F2} escriba:
\excel{=E2-C2}
Copie hasta \texttt{F21}. Una espera de cero indica atención inmediata.
\subsection*{Paso 6. Fin de atención: columna G}
En \texttt{G2} escriba:
\excel{=E2+D2}
Copie hasta \texttt{G21}.
\subsection*{Paso 7. Tiempo total en el sistema: columna H}
En \texttt{H2} escriba:
\excel{=G2-C2}
Copie hasta \texttt{H21}. Este tiempo incluye espera más servicio; no debe confundirse con el tiempo de servicio exclusivamente.
\section{Ejemplo ilustrativo}
Los siguientes valores son una realización posible, no resultados fijos. Cada recálculo genera una simulación diferente.
\begin{center}\small\begin{tabular}{rrrrr}\toprule Cliente & Llegada & Servicio & Espera & Total\\\midrule 1&1,20&2,30&0,00&2,30\\2&2,50&2,70&1,00&3,70\\3&5,80&2,40&0,40&2,80\\4&6,30&2,60&2,30&4,90\\5&10,00&2,50&1,20&3,70\\\bottomrule\end{tabular}\end{center}
\section{Indicadores de la simulación}
En una zona libre de la hoja, registre los siguientes indicadores:
\begin{itemize}[leftmargin=*]
\item Tiempo medio entre llegadas: \texttt{=PROMEDIO(B2:B21)}
\item Tiempo medio de servicio: \texttt{=PROMEDIO(D2:D21)}
\item Tiempo medio de espera: \texttt{=PROMEDIO(F2:F21)}
\item Tiempo medio total en el sistema: \texttt{=PROMEDIO(H2:H21)}
\item Mayor tiempo de espera: \texttt{=MAX(F2:F21)}
\item Número de clientes que esperan: \texttt{=CONTAR.SI(F2:F21;">0")}
\item Proporción de clientes que esperan: \texttt{=CONTAR.SI(F2:F21;">0")/20}
\end{itemize}
\section{Actividades de experimentación}
\begin{enumerate}[leftmargin=*]
\item Complete la tabla para los 20 clientes y verifique que $T_i=W_i+S_i$.
\item Identifique el cliente con la espera máxima y explique por qué ocurrió.
\item Recalcule cinco veces con \texttt{F9} y anote en cada repetición la espera media y la proporción que esperó. Si desea conservar los resultados, cópielos y péguelos como valores.
\item Cambie la media entre llegadas de 3 a 2 minutos (reemplace el factor 3 por 2 en la fórmula de B). Compare los indicadores.
\item Restaure la media entre llegadas a 3 y aumente la media de servicio de 2,5 a 3,5 minutos. Compare los indicadores.
\item Explique por qué los resultados cambian entre simulaciones, incluso sin modificar los parámetros.
\end{enumerate}
\section{Preguntas para discusión}
\begin{enumerate}[leftmargin=*]
\item ¿Qué diferencia existe entre tiempo entre llegadas, instante de llegada y tiempo de espera?
\item ¿Por qué el tiempo de espera no tiene que seguir una distribución exponencial ni normal?
\item ¿Qué efecto tiene la variabilidad del servicio sobre la congestión?
\item ¿Qué ocurre cuando el servicio medio dura más que el intervalo medio entre llegadas?
\item ¿Por qué 20 clientes no son suficientes para obtener estimaciones muy precisas del desempeño a largo plazo?
\end{enumerate}
\section{Alcance del modelo}
El sistema tiene llegadas con intervalos exponenciales y servicios de distribución aproximadamente normal; por ello puede describirse como un modelo de cola \textbf{M/G/1} con disciplina FIFO. Se trata de una simulación didáctica de 20 clientes, no de una validación con datos reales. Para aproximarse a un gemelo digital de una cafetería real sería necesario observar tiempos efectivos, calibrar parámetros y contrastar resultados.
\end{document}
